% ===================================================================
%  TABELO N° 12 — La Staciono. --- La Fervoyi. --- La Navi.
%  Feuillets 82 a 86 du fac-simile = folios 80 a 84.
%  Correspondance : folio imprime = numero de feuillet - 2.
%
%  Les marques « %%K » ne sont pas des commentaires ordinaires : elles
%  portent la CLE D'ALIGNEMENT qui apparie, dans la page de lecture, ce
%  bloc-ci et son homologue francais. La cle se lit
%      t<tableau>-<numero d'alinea numerote>-<rang dans l'alinea>
%  et ne depend d'aucune pagination : les deux editions ne coupent pas
%  leurs pages aux memes endroits, mais elles numerotent les memes
%  alineas — c'est le seul point d'ancrage que les deux livrets
%  partagent, et il est de l'auteur, non du transcripteur.
%
%  NOTA (feuillet 82) : ce tableau ouvre sur un « Noto. » imprime en
%  italique (etiquette « Noto. » en petites capitales sur la page
%  d'ouverture, ailleurs — dans l'appel de bas de page du feuillet 85 —
%  toute l'etiquette est elle-meme en italique) qui n'est ni un alinea
%  numerote ni un intertitre ordinaire : il precede l'alinea « 1. » et
%  fait donc partie de l'apparat d'ouverture. Il est releve comme
%  matiere continue du bloc %%K t12-apar-1, sans cle propre — a
%  l'image de la note de bas de page du tableau 1 (feuillet 9), qui
%  dans le modele texto/io/10-tabelo-01.tex n'a elle non plus aucune
%  cle %%K. Le tableau ne comporte par ailleurs aucune matiere en
%  colonnes ni en horaire tabule : les klok-tabeli qu'il decrit
%  restent de la prose suivie, relevee alinea par alinea.
%
%  Marques d'imprimeur omises : le feuillet 83 (folio 81) porte en bas
%  de page, hors du bloc de texte, la mention de cahier « 6. IDO. » —
%  non relevee, comme le folio courant, qui est deja rendu par
%  \VUfolio.
% ===================================================================

% --- Feuillet 82 (folio 80 : ouverture de tableau) ------------------
\begin{VUpage}[82]{80}
%%K t12-apar-1 apar
\VUsaut{4.0mm}
\VUtitre{15.0pt}{60}{TABELO N\textsuperscript{o} 12}
\VUsaut{2.4mm}
\VUpk{11.4pt}{40}{La Staciono. --- La Fervoyi. --- La Navi.}
\VUsaut{3.4mm}
\VUblancAlinea
\textsc{Noto.} --- \textit{La kloko-tabeli uzata en singla na}\cc
\textit{cioni esas diversa, ni do uzas kom exemplo la}\nl
\textit{kloki di la} \VUgras{Franca tabelo.}

%%K t12-01-1 p
\VUblancAlinea
1. --- Me jus voyajis tra Germania. Ante la\nl
departo, me komprabis \VUgras{klok-libro} \textsuperscript{(15)} por sa\cc
var la departal kloko dil treni. Pos konstatir\nl
ke la maxim oportuna treno esas departonta\nl
de Paris, ye non kloki vespere, ed arivonta\nl
aden Strasburg ye un kloko dek-e-tri minuti,\nl
me preparis mea voyajo, ed en la morga dio,\nl
me igis duktar me a la staciono.

%%K t12-02-1 p
\VUblancAlinea
2. --- Kande me eniris la granda departo\cc
halo, dop olda \VUgras{sacerdoto} \textsuperscript{(2)} rurala, qua manue\nl
kunportis sua \VUgras{voyaj-sako} \textsuperscript{(3)}, multa \VUgras{voya}\cc
\VUgras{jonti} \textsuperscript{(1)} esis ja arivinta.

%%K t12-02-2 p
\VUblancAlinea
Pro ke me volis sendar \VUgras{depesho (telegra}\cc
\VUgras{mo)} \textsuperscript{(5)}, me igis indikar a me da \VUgras{kontrolisto} \textsuperscript{(6)}\nl
\VUgras{la telegraf-kontoro} \textsuperscript{(4)}.

%%K t12-03-1 p
\VUblancAlinea
3. --- Ante pasar aden la \VUgras{varto-salono} \textsuperscript{(7)} me\nl
iris provizar me per \VUgras{bilieto} \textsuperscript{(9)} kun retroveno.

%%K t12-04-1 p
\VUblancAlinea
4. --- Me ne longe vartis an la \VUgras{gicheto} \textsuperscript{(8)},\nl
nam poka personi esis ibe. \VUgras{Soldato} \textsuperscript{(10)}, qua\nl
esis departanta konjede, komplezis cedar a me\nl
sua foyo, tale ke me havis suficanta tempo por\nl
demandar kelk informi de la \VUgras{stacionestro} \textsuperscript{(13)},\nl
qua esis konversanta kun la \VUgras{komisario} \textsuperscript{(12)} dil\nl
surveyado e kun la \VUgras{jendarmo} \textsuperscript{(11)} dejuranta.
\end{VUpage}

% --- Feuillet 83 (folio 81) ------------------------------------------
\begin{VUpage}[83]{81}
%%K t12-tit-1 sub
\VUpk{10.2pt}{40}{Bagaj-enskribado.}
\VUsaut{3.0mm}

%%K t12-05-1 p
\VUblancAlinea
5. --- Komprinte jurnalo che la \VUgras{biblioteko} \textsuperscript{(14)},\nl
me igis ponderar mea \VUgras{bagaji} \textsuperscript{(17)} quin \VUgras{portis}\cc
\VUgras{to} \textsuperscript{(16)} jus adduktis per \VUgras{charioto pakala} \textsuperscript{(19)}; \VUgras{la}\nl
\VUgras{kontoristo} \textsuperscript{(22)} komisita pri la \VUgras{ponder-apara}\cc
\VUgras{to} \textsuperscript{(21)} savigis da me, ke mea kofro pezas tria\cc
dek-e-kin kilogrami; to facis \VUgras{superpezo} de kin\nl
kilogrami; po qua me debas suplementa pago,\nl
ye la \VUgras{gicheto} dil bagaji \textsuperscript{(23)}.

%%K t12-06-1 p
\VUblancAlinea
6. --- Pro ke me esis tro frua, me havis liber\cc
tempo por examenar la cirkulado dil voyajanti\nl
en la staciono. Uli depozis o riprenis sua ne\cc
grava bagaji che la \VUgras{(bagaj)-depozeyo} \textsuperscript{(25)}; altri\nl
konsultis la \VUgras{klok-tabeli} \textsuperscript{(26)}. L'arivanta voya\cc
janti hastis al \VUgras{ekireyo} \textsuperscript{(32)} kun sua \VUgras{valizo} \textsuperscript{(24)} en\nl
manuo. Kelki vartis en la \VUgras{(bagaj)-distribu}\cc
\VUgras{teyo} \textsuperscript{(27)} e prizentis sua \VUgras{bilieto} \textsuperscript{(29)} por ripre\cc
nar sua kofri, \VUgras{kesti} \textsuperscript{(28)} o \VUgras{korbi} \textsuperscript{(30)}.

%%K t12-07-1 p
\VUblancAlinea
7. --- \VUgras{L'oficisti dil acizo} \textsuperscript{(33)} sorgoze exploris\nl
la paki por konstatar kad on havas ulo dekla\cc
renda.

%%K t12-08-1 p
\VUblancAlinea
8. --- Ula voyajanti adiris la \VUgras{bufeto} \textsuperscript{(34)}, ube\nl
\VUgras{garsono} \textsuperscript{(35)} zelis pri servar li. Li mustas has\cc
tar, nam la \VUgras{treno} \textsuperscript{(36)}, qua esas espreso, ne\nl
longe restos en la staciono.

%%K t12-09-1 p
\VUblancAlinea
9. --- Pose me adiris la departo-trotuaro e\nl
serchis direta \VUgras{vagono} \textsuperscript{(37)} por ne esar koaktata\nl
chanjar la vagono dum la voyajo. Me acensis\nl
aden koridor-veturo qua kontenas \VUgras{fako} \textsuperscript{(38)} por\nl
fumeri e fako rezervita por « siorini sola ».
\end{VUpage}

% --- Feuillet 84 (folio 82) ------------------------------------------
\begin{VUpage}[84]{82}
%%K t12-tit-2 sub
\VUpk{10.2pt}{40}{Starto dum la nokto.}
\VUsaut{3.0mm}

%%K t12-10-1 p
\VUblancAlinea
10. --- Pos acensir la \VUgras{ped-apogilo} \textsuperscript{(40)}, klozir\nl
la \VUgras{pordo} \textsuperscript{(39)}, volvir la \VUgras{storo} \textsuperscript{(41)} e lokizir mea\nl
mikra bagaji sur la \VUgras{reto} \textsuperscript{(43)}, me sidis sur la\nl
\VUgras{benketo} \textsuperscript{(42)}. Vidante la \VUgras{lampisto} \textsuperscript{(44)} acendar la\nl
lampi, me pensis, ke me devos pasar nokto en\nl
la vagono e me advokis l'oficisto komisita pri\nl
la lokaco dil \VUgras{kapkuseni} \textsuperscript{(45)} por demandar un\nl
de oli.

%%K t12-11-1 p
\VUblancAlinea
11. --- La tren-konduktoro venis verifikar la\nl
bilieti. Me questionis lu pro quo ni ne ja de\cc
partis, nam esas la justa kloko. Il respondis\nl
a me ke la \VUgras{trenestro} \textsuperscript{(46)} okupesas per pozigar\nl
bicikli en la \VUgras{furgono} \textsuperscript{(47)}. Pluse, on ne ja chan\cc
jabis omna pedvarmigili \textsuperscript{(48)}.

%%K t12-12-1 p
\VUblancAlinea
12. --- Ma ni esis balde departonta, nam la\nl
\VUgras{fairisto} \textsuperscript{(50)}, sur sua \VUgras{tendro} \textsuperscript{(49)} prenis karbono\nl
por ecitar la fairo di la \VUgras{lokomotivo} \textsuperscript{(54)}, e\nl
\VUgras{la mashinisto} \textsuperscript{(51)} vartis nur la signalo di la\nl
\VUgras{staciono-subchefo} \textsuperscript{(52)}, qua esis sur la \VUgras{tro}\cc
\VUgras{tuaro} \textsuperscript{(53)}.

%%K t12-13-1 p
\VUblancAlinea
13. --- La \VUgras{kaldiero} \textsuperscript{(56)} di la lokomotivo ma\cc
tide grondis; la \VUgras{kameno} \textsuperscript{(55)} vomis torenti de\nl
fumuro mixita kun \VUgras{vaporo} \textsuperscript{(60)}. Stridanta \VUgras{si}\cc
\VUgras{flo} \textsuperscript{(59)} sonegis e la \VUgras{pistoni} \textsuperscript{(58)} moveskis, impul\cc
sante la \VUgras{roti} \textsuperscript{(57)} di la pezoza mashino.

%%K t12-tit-3 sub
\VUsaut{3.0mm}
\VUpk{10.2pt}{40}{Departo !}
\VUsaut{3.0mm}

%%K t12-14-1 p
\VUblancAlinea
14. --- La rapideso dil treno, kuranta sur la\nl
\VUgras{reli} \textsuperscript{(64)} esis acelerata; la \VUgras{trotuari} \textsuperscript{(65)} desaparis;\nl
ni preterpasis la \VUgras{latrini} \textsuperscript{(66)}, ed anke vagonaro
\parplein
\end{VUpage}

% --- Feuillet 85 (folio 83) ------------------------------------------
\begin{VUpage}[85]{83}
%%K t12-noto-1 noto
\VUnotes{143.2mm}{%
\textit{(*) Noto.} --- \textit{Komprenende la voyajo-deskripto esas segun}\nl
\textit{la pre-milita frontiero.}%
}

%%K t12-14-1 p suite
\VUcontinue
garita sur l'altra \VUgras{voyo} \textsuperscript{(63)} : \VUgras{dormo-vagoni} \textsuperscript{(67)},\nl
\VUgras{restor-vagono} \textsuperscript{(68)}, \VUgras{posto-vagono} \textsuperscript{(69)}.

%%K t12-15-1 p
\VUblancAlinea
15. --- Pose, la \VUgras{staciono dil vari} \textsuperscript{(71)} pasis\nl
avan nia okuli. Sub la hangari, oficisti kargis\nl
la \VUgras{vari} \textsuperscript{(72)} expediita per nerapida treno. \VUgras{Es}\cc
\VUgras{quadani} \textsuperscript{(76)} proxim l'\VUgras{aquo-rezervuyo} \textsuperscript{(73)} ma\cc
novris \VUgras{bestio-vagoni} \textsuperscript{(75)} e \VUgras{platforma vagoni} \textsuperscript{(79)}\nl
por formacar \VUgras{varo-treno} \textsuperscript{(80)}.

%%K t12-16-1 p
\VUblancAlinea
16. --- Ni transiris la lasta \VUgras{bifurkeyo} \textsuperscript{(78)},\nl
proxim qua stacis l'agulisto; ni preteriris la\nl
\VUgras{signal-disko} \textsuperscript{(86)} e la staciono fore desaparis.

%%K t12-17-1 p
\VUblancAlinea
17. --- Balde ni trairis \VUgras{fera ponto} \textsuperscript{(74)} qua\nl
transpasas la \VUgras{fluvio} \textsuperscript{(81)}. Pos trairir ta ponto,\nl
ni alongiris la rivero, sur qua \VUgras{remorkeri} \textsuperscript{(82)}\nl
potenta tiris pezoza \VUgras{kargo-bateli} \textsuperscript{(83)}. Ni pluse\nl
vidis \VUgras{voyaj-navo} \textsuperscript{(84)}, kande ol esis abordanta\nl
\VUgras{pontono} \textsuperscript{(85)}.

%%K t12-18-1 p
\VUblancAlinea
18. --- Pos rapidege preterpasir plura sta\cc
cioni, ni trairis kelka \VUgras{tuneli} \textsuperscript{(87)} e lasis dop ni\nl
\VUgras{dometi} sennombra \textsuperscript{(88)} di \VUgras{baril-gardisti.} Ber\cc
sate dal sukusado dil treno, me profunde dor\cc
meskis.

%%K t12-tit-4 sub
\VUsaut{3.0mm}
\VUpk{10.2pt}{40}{Deutsch-Avricourt-Staciono.}
\VUsaut{3.0mm}

%%K t12-19-1 p
\VUblancAlinea
19. --- Me esis subite vekigata dal voco di\nl
oficisto krianta : « Deutsch-Avricourt ! \textsuperscript{(*)}.\nl
Omni decensez ! » Eventis l'inspekto da la do\cc
gano ed omna tedanta formalaji di la frontiero.\nl
Me mustis konformigar mea posh-horlojo a la\nl
\VUgras{horlojego} \textsuperscript{(89)} di la staciono, qua esis tro frua
\parplein
\end{VUpage}

% --- Feuillet 86 (folio 84) ------------------------------------------
\begin{VUpage}[86]{84}
%%K t12-19-1 p suite
\VUcontinue
per kinadek e kin minuti; apene vekigita, me\nl
desfacile distingis la \VUgras{grand indexo} \textsuperscript{(90)} de la \VUgras{mi}\cc
\VUgras{kra} \textsuperscript{(91)}; la \VUgras{cifro-plako} \textsuperscript{(92)} ipsa esis tote kon\cc
fuza pro la matinal nebuleto.

%%K t12-20-1 p
\VUblancAlinea
20. --- Dum ke la doganisti facis sua ins\cc
pekto, me ocitante dechifris la \VUgras{afishi} \textsuperscript{(93)} qui\nl
dekoras la stacion-muri. Me dejunis por pasi\cc
gar la tempo , konsultis la mapo dil Germana\nl
\VUgras{reto} \textsuperscript{(94)} por vidar quanta disto ankore separas\nl
me de Strasburg, e me rieniris mea vagono,\nl
rifortigita dal espero balde atingor la skopo di\nl
mea voyajo.

\VUsaut{6.0mm}
\VUfilet{22mm}
\end{VUpage}
